% =====================================================================
%  Section 8.5: Long-time behaviour, the N-wave.  Optional: remove the
%  \input line in main.tex to drop this subsection.  The corrected
%  constants P, Q were checked numerically against the Lax--Oleinik
%  formula for random piecewise-constant data (L^1 distance to the N-wave
%  halves each time t is multiplied by 4, i.e. O(t^{-1/2})).
% =====================================================================

\subsection{Long-time behaviour: the N-wave}
\label{subsec:n-wave}

This final extension describes the universal shape that entropy
solutions with compactly supported bounded initial data approach as $t \to
\infty$, regardless of the detailed shape of $u_0$. The point to keep in
mind is that the two constants specifying this shape are \emph{not} the
total positive and negative masses of $u_0$: positive and negative parts
of the data can annihilate each other across shocks.

\subsubsection{An exact self-similar solution}
\label{subsubsec:n-wave-exact}

\begin{proposition}[The N-wave]
\label{prop:n-wave-exact}
Fix constants $P, Q > 0$. The function
\begin{equation}
  \label{eq:n-wave}
  \mathcal{N}(x,t) = \begin{cases}
    0, & x \le -\sqrt{2Qt}, \\
    x/t, & -\sqrt{2Qt} < x < \sqrt{2Pt}, \\
    0, & x \ge \sqrt{2Pt},
  \end{cases}
  \qquad t>0,
\end{equation}
is an entropy solution of $u_t+u\,u_x=0$ on $\mathbb{R}\times[t_0,\infty)$
for every $t_0>0$ (with data $\mathcal{N}(\cdot,t_0)$). Moreover
$\int \mathcal{N}^+(x,t)\,dx=P$, $\int \mathcal{N}^-(x,t)\,dx=Q$ for all $t>0$.
\end{proposition}

\begin{proof}
Apply Proposition~\ref{prop:piecewise-weak} on $t\ge t_0$ with the two
curves $x_-(t)=-\sqrt{2Qt}$ and $x_+(t)=\sqrt{2Pt}$, which are $C^1$ there.
(i) On the middle interval $\mathcal{N}=x/t$ is the self-similar fan of
Section~\ref{subsec:rarefaction-general}, which solves the PDE
classically; the outer regions are constant. (ii) At the right edge the
one-sided limits are $u_\ell = x_+(t)/t = \sqrt{2P/t}$ and $u_r=0$; since
$\dot x_+(t) = \sqrt{P/(2t)} = \tfrac12\big(\sqrt{2P/t}+0\big)$, this is
exactly the Rankine--Hugoniot speed \eqref{eq:rankine-hugoniot-burgers}.
At the left edge $u_\ell=0$, $u_r=-\sqrt{2Q/t}$, and $\dot x_-(t) =
-\sqrt{Q/(2t)} = \tfrac12\big(0+(-\sqrt{2Q/t})\big)$, again matching
\eqref{eq:rankine-hugoniot-burgers}. (iii) The initial trace at $t_0$ is
$\mathcal{N}(\cdot,t_0)$ by definition. Hence $\mathcal{N}$ is a weak solution there.

For the entropy condition \eqref{eq:oleinik}: both jumps are downward
($u_\ell>u_r$), and inside the middle interval $N_x=1/t$. If $x$ and $x+z$
both lie in the middle interval, $\mathcal{N}(x+z,t)-N(x,t)=z/t$; if $x$ lies left of
it, $\mathcal{N}(x+z,t)\le z/t$ because $x<0$. The only remaining case with a
possible increase is $x<0$ inside the middle interval and $x+z$ beyond the
right edge; then $\mathcal{N}(x+z,t)-N(x,t)=-x/t=|x|/t\le z/t$, because
$x+z>0$ forces $|x|<z$. So \eqref{eq:oleinik} holds. Finally $\int_0^{\sqrt{2Pt}}(x/t)\,dx=P$
and $\int_{-\sqrt{2Qt}}^0 (-x/t)\,dx=Q$.
\end{proof}

The profile \eqref{eq:n-wave} is called an \emph{N-wave}: it rises
linearly through the origin and is bounded by a shock on each side, so
that its graph traces the shape of the letter N (reversed), decaying in
amplitude like $t^{-1/2}$ and spreading in width like $t^{1/2}$ as
$t\to\infty$. (If $P=0$ or $Q=0$ the corresponding shock is absent and
the profile is a single triangle ending in one shock; $P=Q=0$ gives
$\mathcal{N}\equiv0$.) The function $\mathcal{N}$ is unbounded as $t\to0^+$, so it is not a
weak solution on $t>0$ in the sense of Definition~\ref{def:weak-solution};
as $t\to0^+$ it converges, in the sense of distributions, to $(P-Q)\delta_0$:
both edges collapse to the origin, a state concentrated entirely at a
point.

\subsubsection{Shock merging and cancellation}
\label{subsubsec:shock-merging}

Two mechanisms are visible as general solutions evolve towards the
N-wave; both are illustrations of why such a limit is plausible, and
neither, by itself, proves the convergence theorem of
Section~\ref{subsubsec:n-wave-convergence}. The first is that colliding
shocks merge into one, governed by exactly the Riemann-problem machinery
already developed.

\begin{lemma}[Shock merging]
\label{lem:shock-merging}
Let $u_1>u_2>u_3$, and consider initial data equal to $u_1$ for $x<a$,
$u_2$ for $a<x<b$, and $u_3$ for $x>b$. For small $t>0$ the entropy
solution consists of two shocks moving at speeds $s_{12}=(u_1+u_2)/2$
and $s_{23}=(u_2+u_3)/2$; since $s_{12}>s_{23}$, they collide at the
finite time $t^\times = 2(b-a)/(u_1-u_3)$, after which the entropy
solution is the single shock connecting $u_1$ to $u_3$ at speed
$s_{13}=(u_1+u_3)/2$.
\end{lemma}

\begin{proof}[Proof sketch]
Each jump individually satisfies the Lax condition since
$u_1>u_2>u_3$, so Theorem~\ref{thm:riemann-solution}(a) gives the two
shock speeds above. Since $s_{12}-s_{23}=\tfrac12(u_1-u_3)>0$, the shock
positions $a+s_{12}t$ and $b+s_{23}t$ meet at the stated finite
$t^\times$. At that instant the middle state has shrunk to zero width,
leaving the two-state data $u_1,u_3$ on either side of the collision
point; since $u_1>u_3$, Theorem~\ref{thm:riemann-solution}(a) applied
afresh from time $t^\times$ gives the merged shock at speed $s_{13}$.
\end{proof}

\begin{remark}[Scope of the lemma]
Lemma~\ref{lem:shock-merging} covers only shock--shock collisions. In
general a shock also absorbs rarefaction fans, and a shock separating a
positive state (on its left) from a negative one (on its right) cancels
positive against negative mass. The following example shows this
cancellation concretely and why it matters.
\end{remark}

\begin{example}[Cancellation]
\label{ex:cancellation}
Let $u_0=+1$ on $(0,1)$, $u_0=-1$ on $(1,2)$, and $u_0=0$ elsewhere. For
$0<t<1$ the entropy solution is
\begin{equation*}
  u=\begin{cases}
    0,&x<0,\\ x/t,&0<x<t,\\ 1,&t<x<1,\\ -1,&1<x<2-t,\\
    (x-2)/t,&2-t<x<2,\\ 0,&x>2,
  \end{cases}
\end{equation*}
(two fans, and a stationary shock at $x=1$, since $u_\ell=1>u_r=-1$ and
$(u_\ell+u_r)/2=0$), as one checks with
Proposition~\ref{prop:piecewise-weak}. At $t=1$ the fans reach the shock,
and for $t\ge1$
\begin{equation*}
  u(x,t)=\begin{cases}
    x/t,&0<x<1,\\ (x-2)/t,&1<x<2,\\ 0,&\text{otherwise,}
  \end{cases}
\end{equation*}
with the stationary shock at $x=1$ ($u_\ell=1/t$, $u_r=-1/t$) and
continuity elsewhere; it satisfies \eqref{eq:oleinik} because its only
jump is downward and $u_x=1/t$ elsewhere. Hence for $t\ge1$,
$\|u(\cdot,t)\|_{L^1}=\frac1{2t}+\frac1{2t}=\frac1t\to0$, even though
$\int u_0^+=\int u_0^-=1$.
\end{example}

\subsubsection{Convergence of general solutions to the N-wave}
\label{subsubsec:n-wave-convergence}

The correct constants are read off from the primitive of $u_0$, not from
its positive and negative parts.

\noindent The following is a \emph{cited} theorem: it is stated here without
proof, and the reader should consult \cite{lax1957} and
\cite[\S3.4.5]{evans2010} for the precise hypotheses, norm and proof.

\begin{theorem}[{Asymptotic N-wave decay; Lax \cite{lax1957}, see also \cite[\S3.4.5]{evans2010}}]
\label{thm:n-wave-decay}
Let $u_0 \in L^\infty(\mathbb{R})$ have compact support, and set
\begin{equation}
  \label{eq:pq-areas}
  V(y)=\int_{-\infty}^y u_0(x)\,dx,\qquad M=\int_{\mathbb{R}}u_0\,dx,\qquad
  Q=-\min_y V(y),\qquad P=M-\min_y V(y).
\end{equation}
Then $P,Q\ge0$, $P-Q=M$, and the entropy solution $u$ of \eqref{eq:ivp}
satisfies
\begin{equation}
  \|u(\cdot,t) - \mathcal{N}(\cdot,t)\|_{L^1(\mathbb{R})} \to 0 \qquad \text{as } t \to \infty,
\end{equation}
where $\mathcal{N}$ is the N-wave \eqref{eq:n-wave} with these constants $P,Q$.
\end{theorem}

% TODO: insert the exact theorem number from Evans Section 3.4.5 and the
% relevant theorem/section of Lax (1957) in the references.

\noindent (The primitive $V$ here is based at $-\infty$, whereas in
Remark~\ref{rem:lax-oleinik} it is based at $0$; the two differ by a constant,
which changes neither the minimisers of $V(y)+(x-y)^2/2t$ nor $u$.)
Equivalently $Q=\max_y\int_{-\infty}^y(-u_0)$ and
$P=\max_y\int_y^\infty u_0$, so in general
$P\le\int u_0^+$ and $Q\le\int u_0^-$, with equality exactly when $u_0\le0$
to the left and $u_0\ge0$ to the right of some point (a ``pre-formed''
N-shape, with no cancellation). Example~\ref{ex:cancellation} has $V\ge0$,
so $Q=0$ and $P=M-0=0$: the limiting N-wave is $\mathcal{N}\equiv0$, consistent with
$\|u\|_{L^1}=1/t\to0$.

\paragraph{Heuristic derivation of the constants (not a proof of
convergence).} The shock-merging lemma and the computation below
\emph{motivate} the form of the theorem and the formula for $P$ and $Q$;
they do not establish the $L^1$ convergence, which requires the estimates
in the cited sources. The constants \eqref{eq:pq-areas} arise from the
Lax--Oleinik formula \eqref{eq:lax-oleinik} of
Remark~\ref{rem:lax-oleinik} by a short computation. Here $V(y)=0$ to the
left of the support, $V(y)=M$ to its right, and $\min V=-Q$ is attained at
some bounded point $y_0$. For $x$ of order $\sqrt t$ the minimiser of
$V(y)+(x-y)^2/(2t)$ is one of two candidates:
\begin{itemize}
  \item $y\approx x$, with value $0$ for $x$ left of the support and $M$
    for $x$ right of it, giving $u=(x-y)/t=0$;
  \item $y\approx y_0$, with value $-Q+(x-y_0)^2/(2t)\approx -Q+x^2/(2t)$,
    giving $u=(x-y_0)/t\approx x/t$.
\end{itemize}
The second beats the first on the left when $-Q+x^2/(2t)<0$, i.e.\
$x>-\sqrt{2Qt}$, and on the right when $-Q+x^2/(2t)<M$, i.e.\
$x<\sqrt{2(M+Q)t}=\sqrt{2Pt}$. This produces exactly the profile
\eqref{eq:n-wave}, with shocks (ties between the two candidates) at
$x=-\sqrt{2Qt}$ and $x=\sqrt{2Pt}$. Making this rigorous --- controlling
the error in $L^1$ --- is the content of the cited theorem, and is not
attempted here. In this
picture the conserved quantity is only $\int u\,dx=M=P-Q$; the
individual masses $P$ and $Q$ are \emph{not} conserved, since positive and
negative parts cancel at shocks as in Example~\ref{ex:cancellation}.

\subsubsection{A numerical check of the constants}
\label{subsubsec:n-wave-numerics}

The following is an \emph{observation}, not a proof. For piecewise-constant
compactly supported $u_0$ the entropy solution can be computed (to grid
resolution) directly from the Lax--Oleinik formula
\eqref{eq:lax-oleinik}, with $V$ the exact primitive
(\texttt{code/nwave\_check.py}). We compare it in $L^1$ with the N-wave built
from the corrected constants \eqref{eq:pq-areas} and, for contrast, from the
\emph{naive} constants $P=\int u_0^+$, $Q=\int u_0^-$. Three data sets were
used: (a) $-1$ on $(0,1)$, $+1$ on $(1,2)$ (no cancellation); (b)
$+1$ on $(0,1)$, $-1$ on $(1,2)$ (Example~\ref{ex:cancellation}); (c) six
alternating-sign pieces with values $1.2,-0.8,1.6,-1.4,0.9,-0.6$ on
$(0,0.7,1.5,2.1,3.0,3.8,4.4)$.

\begin{table}[h]
\centering
\begin{tabular}{l c c c c c}
  data & corrected $(P,Q)$ & naive $(P,Q)$ & $L^1$ error at $t=6400$ & naive error & fitted slope \\\hline
  (a) & $(1,1)$ & $(1,1)$ & $7.0\times10^{-2}$ & $7.0\times10^{-2}$ & $-0.49$ \\
  (b) & $(0,0)$ & $(1,1)$ & $1.6\times10^{-4}$ & $2.0$ & $-1.00$ \\
  (c) & $(0.36,0.10)$ & $(2.52,2.26)$ & $9.4\times10^{-2}$ & $4.4$ & $-0.44$
\end{tabular}
\caption{$L^1$ distance between the entropy solution and the N-wave, for
$t\in\{25,100,400,1600,6400\}$; ``fitted slope'' is the log--log slope of
the corrected-constant error over $t\ge100$.}
\label{tab:nwave}
\end{table}

With the corrected constants the error decays in all three cases. In (a)
the slope is $-0.49$, consistent with $t^{-1/2}$; in (b) the limiting
N-wave is $\mathcal{N}\equiv0$ and the error equals $\|u(\cdot,t)\|_{L^1}=1/t$ exactly
(slope $-1.00$, as computed in Example~\ref{ex:cancellation}); in (c) the
slope over this range is $-0.44$ and still changing, so no rate is claimed
there. With the naive constants the error does not decay in (b) or (c): it
stays near $\|\mathcal{N}_{\rm naive}-\mathcal{N}_{\rm corr}\|_{L^1}
=(P_{\rm n}-P_{\rm c})+(Q_{\rm n}-Q_{\rm c})$ (the extra mass of the larger
naive wave, which has the same profile on the common interval), namely $2.0$
in (b) and $2.16+2.16=4.32$ in (c) (observed $4.35$ at $t=6400$). The naive
profile itself has $L^1$ norm $P+Q=2.0$ and $4.78$ respectively; in (b) the
two quantities coincide, which can hide the distinction. Either way, the
positive and negative masses of $u_0$ are not the correct constants.
